(1) 最小耦合的环向动量为 $-i\hbar R^{-1}\partial_\theta-qA_\theta$，故 $H=(2mR^2)^{-1}[-i\hbar\partial_\theta-q\Phi/(2\pi)]^2$。环上虽然 $B=0$，矢势的闭合线积分 $\oint\mathbf A\cdot d\mathbf l=\Phi$ 不为零，不能用单值规范变换在所有磁通下同时消除。
(2) 单值周期性 $\psi(\theta+2\pi)=\psi(\theta)$ 给归一化本征态 $\psi_n=e^{in\theta}/\sqrt{2\pi}$，$n\in\mathbb Z$。相应能谱 $E_n(\Phi)=\hbar^2[n-\Phi/\Phi_0]^2/(2mR^2)$，其中 $\Phi_0=h/q=2\pi\hbar/q$。
(3) 当 $\Phi\to\Phi+\Phi_0$，有 $E_n(\Phi+\Phi_0)=E_{n-1}(\Phi)$，所以所有能级组成的集合与热力学谱量对一个磁通量子周期性，而单个固定标签的抛物线会换标。在 $\Phi=\Phi_0/2$ 时 $n=0$ 和 $n=1$ 两态都有 $E=\hbar^2/(8mR^2)$，出现简并。
(4) 对磁通求导，$I_n=-\partial E_n/\partial\Phi=\hbar^2(n-\Phi/\Phi_0)/(mR^2\Phi_0)$，在所用正电荷方向约定下符号随 $n-\Phi/\Phi_0$ 改变。磁通通过环的拓扑相位改变周期边界下的动量量子化，即使粒子轨迹处 $B=0$，谱与持续电流仍可随磁通变化。
